% Part: lambda-calculus
% Chapter: introduction
% Section: basic-pr-lambda

\documentclass[../../../include/open-logic-section]{subfiles}

\begin{document}

\olfileid{lam}{int}{bas}
\olsection{मूलभूत आदिम पुनरावर्ती फलने \usetoken{S}{lambda definable} असतात}

\begin{lem}
$\Zero$, $\Succ$ आणि $\Proj{n}{i}$ ही फलने
!!{lambda definable} आहेत.
\end{lem}

\begin{proof}
%READERNOTE{SOL6-A108: आधीच्या recursive-function व्याख्येतील शून्य फलन एकस्थानी आहे. चर्चचा शून्य अंक आणि तो अंक नेहमी परत करणारे एकस्थानी फलन यांत अतिरिक्त अमूर्तीकरणाचा फरक आहे.}
येथे $\Zero$ हे आधी परिभाषित केलेले एकस्थानी शून्य फलन आहे.
त्याला दर्शवणारे पद
$\lambd[w][\lambd[x][\lambd[y][y]]]$
असे घ्या. बाहेरचे अमूर्तीकरण निविष्ट बाजूला ठेवते आणि
चर्चचा शून्य अंक परत करते.

उत्तरवर्ती फलन~$\Succ$ ची
$\fn{Succ}(u) = \lambd[x][\lambd[y][x(uxy)]]$ अशी
व्याख्या करतात. हे कसे चालते याचा विचार करा: पुनरावर्तक
मानलेल्या प्रत्येक अंकासाठी $\num{n}$ आणि प्रत्येक
फलनासाठी~$f$, $\fn{Succ}(\num{n},f)$ हे असे फलन आहे
की निविष्ट~$y$ वर ते~$y$ पासून सुरू करून $f$ चे $n$
वेळा उपयोजन करते आणि त्यानंतर आणखी एकदा उपयोजन करते.

प्रक्षेपणांबद्दल अधिक काही सांगण्यासारखे नाही:
$\fn{Proj}^n_i(x_0, \dots,
x_{n-1}) = x_i$. दुसऱ्या शब्दांत, आपल्या रूढींनुसार,
$\fn{Proj}^n_i$ हे लॅम्डा पद
$\lambd[x_0][\dots \lambd[x_{n-1}][x_i]]$ आहे.
\end{proof}

\end{document}
